\phantomsection
\addchaptertocentry{Удиртгал}
\label{Introduction}

\section*{Удиртгал}

Сүүлийн жилүүдэд гадаадын их, дээд сургуульд суралцах сонирхолтой ахлах ангийн сурагчдын тоо нэмэгдэж, элсэлтийн бэлтгэл ажлыг илүү зохион байгуулалттайгаар удирдах хэрэгцээ өсөж байна. Гадаадын их сургуульд өргөдөл гаргах үйл явц нь олон төрлийн материал бүрдүүлэх, хугацаа хянах, багш болон сурагчийн хооронд тогтмол харилцах зэрэг олон шат дамжлагатай байдаг. Гэвч одоогийн нөхцөлд эдгээр үйл ажиллагаа нь и-мэйл, чат, Excel файл, цаасан баримт зэрэг салангид хэрэгслүүд ашиглан явагддаг тул мэдээлэл алдагдах, хугацаа хоцрох, даалгаврын явц хянахад хүндрэл үүсэх зэрэг асуудлууд гардаг.

Энэхүү систем нь сургуулийн зөвлөх багш, сурагч болон сургуулийн удирдлагад зориулсан вебд суурилсан нэгдсэн платформ юм. Системийн хүрээнд зөвлөх багш сурагчдад даалгавар болон дэд даалгавар үүсгэх, гүйцэтгэлийн явцыг хянах, файл солилцох, бодит цагийн чат ашиглан харилцах боломжтой. Сурагчид өөрийн даалгаврын төлөвийг харах, материал илгээх, мэдэгдэл хүлээн авах, тодорхойлох захидлын хүсэлт илгээх зэрэг үйлдлийг хийх боломжтой. Мөн сургуулийн удирдлагын түвшинд элсэлтийн статистик болон тайлангийн мэдээллийг нэг дороос харах боломжийг бүрдүүлнэ.

Энэхүү системийг хөгжүүлснээр зөвлөх багшийн ажлын ачааллыг бууруулах, сурагчдын гүйцэтгэлийн хяналтыг сайжруулах, элсэлтийн үйл явцыг илүү ил тод, зохион байгуулалттай болгох нөхцөл бүрдэнэ.
\subsection*{Зорилго}
Энэхүү төслийн зорилго нь ахлах сургуулийн сурагчдын гадаадын их, дээд сургуульд элсэх бэлтгэл үйл явцыг удирдах, зөвлөх багш, сурагч болон сургуулийн удирдлагын харилцаа, даалгаврын менежмент, баримт бичгийн солилцоог нэгтгэсэн вебд суурилсан системийг хөгжүүлэхэд оршино.

\subsection*{Зорилт}
Уг зорилгод хүрэхийн тулд дараах зорилтуудыг тодорхойлсон. Үүнд:
\begin{enumerate}[nosep]
    \item Ахлах ангийн сурагчид болон зөвлөх багш нарын элсэлтийн бэлтгэлтэй холбоотой үйл ажиллагаа, тулгамдаж буй асуудлуудыг судлах;
    \item Гадаадын их, дээд сургуульд элсэх үйл явцын үе шат, шаардлагыг тодорхойлж, системийн функциональ болон функциональ бус шаардлагуудыг боловсруулах;
    \item Даалгавар болон дэд даалгаврын удирдлага, файл солилцоо, чат харилцаа, мэдэгдэл, тайлангийн үндсэн модулиудыг агуулсан веб системийг Next.js технологид суурилан хөгжүүлэх;
    \item Системийн өгөгдлийн сан, архитектур болон UML загваруудыг боловсруулж, хэрэгжүүлэлтийн бүтэцтэй уялдуулах;
    \item Веб технологид суурилан хэрэгжүүлэх;
    \item Хөгжүүлсэн системийг бодит хэрэглэгчдийн оролцоотой туршиж, системийн ашиглахад хялбар байдал болон гүйцэтгэлд дүн шинжилгээ хийх.
\end{enumerate}

\subsection*{Хамрах хүрээ}
Энэхүү төсөл нь дараах хүрээг хамарна:
\begin{itemize}
    \item Ахлах сургуулийн зөвлөхийн үйл ажиллагаанд зориулсан веб платформ хөгжүүлэх;
    \item Сурагчдад даалгавар болон дэд даалгавар үүсгэх, гүйцэтгэлийг хянах модуль;
    \item Файл хавсаргах, баримт бичиг хадгалах болон солилцох модуль;
    \item Бодит цагийн чат болон мэдэгдлийн систем;
    \item Тодорхойлох захидлын хүсэлт удирдах модуль;
    \item Элсэлтийн үр дүн болон карьер тестийн мэдээлэл хадгалах модуль;
    \item Сургуулийн удирдлагад зориулсан тайлан, статистикийн мэдээллийн хэсэг;
    \item Олон хэрэглэгчийн эрхийн түвшинд ажиллах хандалтын удирдлагын систем.
\end{itemize}